\chapter{Evaluation Framework}
\label{chap:evaluation}

This section specifies the evaluation components required for Simula methodology
compliance, as described in Sections 2.3, 3.1, and 3.3 of the research paper.

\section{Calibrated Complexity Scoring (Algorithm 3)}

Raw per-sample complexity scores are noisy due to model overconfidence.
The complexity calibrator uses batch-wise relative scoring with Elo ranking.

\subsection{Elo Rating Algorithm}

The Elo update follows the standard formula (Elo \& Sloan, 1978), adapted for
complexity comparisons. Given two examples $a$ and $b$ with ratings $R_a$ and
$R_b$, where $a$ has higher raw complexity score:

\begin{enumerate}[nosep]
    \item Compute expected outcome:
          $E_a = \frac{1}{1 + 10^{(R_b - R_a)/400}}$
    \item Update ratings with K-factor = 32:
          $R'_a = R_a + K \cdot (1 - E_a)$,
          $R'_b = R_b + K \cdot (0 - (1 - E_a))$
\end{enumerate}

\begin{lstlisting}[language=Python,caption={Elo computation (simula\_complexity\_calibrator.py)}]
def _update_elo(self, id_a: str, id_b: str, outcome_a: float) -> None:
    """
    Update Elo ratings for a pairwise comparison.
    
    Args:
        id_a, id_b: Example identifiers
        outcome_a: 1.0 if A wins (higher complexity), 0.0 if B wins, 0.5 for draw
        
    Standard Elo formula per Elo & Sloan (1978).
    """
    rating_a = self._scores[id_a].elo_rating
    rating_b = self._scores[id_b].elo_rating
    
    # Expected outcomes
    expected_a = 1 / (1 + math.pow(10, (rating_b - rating_a) / 400))
    expected_b = 1 - expected_a
    
    # K-factor = 32 (standard for chess, appropriate for complexity)
    K = self.config.elo_k_factor  # default 32
    
    # Update ratings
    self._scores[id_a].elo_rating = rating_a + K * (outcome_a - expected_a)
    self._scores[id_b].elo_rating = rating_b + K * ((1 - outcome_a) - expected_b)
\end{lstlisting}

\subsection{Usage}

\begin{lstlisting}[language=Python,caption={Complexity calibrator usage}]
from src.training.pipeline import (
    SimulaComplexityCalibrator,
    ComplexityConfig,
)

calibrator = SimulaComplexityCalibrator(
    llm_client=llm,
    config=ComplexityConfig(
        batch_size=5,   # BS=5 from Figure 10
        n_samples=5,    # N=5 from Figure 10
        enable_elo=True,
        elo_k_factor=32,  # Standard K-factor
    ),
)

# Calibrate all examples
scores = await calibrator.calibrate(examples)

# Split by complexity (Figure 7)
low, mid, high = calibrator.split_by_complexity(
    examples,
    low_percentile=0.4,
    high_percentile=0.6,
)

# Export scores
calibrator.export_scores("complexity_elo_scores.json")
\end{lstlisting}

\textbf{Key insight from paper:} High-complexity data yields +10\% accuracy on GSM8k
at 64k samples (Figure 7), but can \emph{hurt} performance when the teacher model
is weak (LEXam case).

\section{Taxonomic Coverage Evaluation}

The coverage evaluator measures Level Ratio Coverage---the proportion of unique
taxonomy nodes covered at each level:

\begin{lstlisting}[language=Python,caption={Coverage evaluator usage}]
from src.training.pipeline import (
    SimulaCoverageEvaluator,
    CoverageConfig,
    identify_coverage_gaps,
)

evaluator = SimulaCoverageEvaluator(
    llm_client=llm,
    config=CoverageConfig(
        max_level=4,
        sample_size=1000,
    ),
)

# Evaluate against all taxonomies
reports = await evaluator.evaluate_multiple(examples, taxonomies)

# Identify gaps
gaps = identify_coverage_gaps(reports, threshold=0.5)
print(f"Under-covered nodes: {gaps}")

# Export
evaluator.export_all_reports(reports, "coverage_report.json")
\end{lstlisting}

\section{Embedding-Based Diversity Metrics}

Following Yu et al. (2023), diversity is measured via cosine distance in embedding
space. \textbf{Global diversity} is the average pairwise distance across the dataset.
\textbf{Local diversity} is the average k-NN distance.

\begin{lstlisting}[language=Python,caption={Diversity analyzer usage}]
from src.training.pipeline import (
    SimulaDiversityAnalyzer,
    DiversityConfig,
)

analyzer = SimulaDiversityAnalyzer(
    config=DiversityConfig(
        embedding_model="sentence-transformers/all-MiniLM-L6-v2",
        knn_k=10,
    ),
)

report = analyzer.analyze(examples, text_field="question")
print(f"Global: {report.global_diversity:.4f}")
print(f"Local:  {report.local_diversity:.4f}")

# Identify outliers (near-duplicates or errors)
low_ids, high_ids = analyzer.identify_outliers(report)
\end{lstlisting}

\textbf{Key insight:} Global and Local diversification have \emph{additive} effects
(Figure 4). Optimizing only one is insufficient.

\section{Critic Calibration Validation}

The critic validator measures $p(y)$ and $p(y_{\text{corrupt}})$ to ensure the
double-critic is properly calibrated:

\begin{lstlisting}[language=Python,caption={Critic validator usage}]
from src.training.pipeline import (
    SimulaCriticValidator,
    CriticValidationConfig,
)

validator = SimulaCriticValidator(
    llm_client=llm,
    config=CriticValidationConfig(
        validation_split=0.05,
        max_samples=200,
    ),
)

report = await validator.validate(examples, schema_context)
print(f"p(y)        = {report.p_y:.3f}")
print(f"p(y_corrupt)= {report.p_y_corrupt:.3f}")
print(f"Calibrated: {report.is_well_calibrated}")
\end{lstlisting}

\textbf{Key insight:} Critic effectiveness degrades with task complexity (Figure 3b).
For effective rejection sampling, we need $p(y) > 0.8$ and $p(y_{\text{corrupt}}) < 0.3$.

\section{Student-Teacher Gap Analysis}

The paper shows that downstream performance saturates when the student bridges
approximately 83\% of the teacher-student gap (Section 4.3):

\begin{lstlisting}[language=Python,caption={Gap analysis}]
from src.training.pipeline import SimulaConfig

config = SimulaConfig()
config.training.teacher_accuracy = 0.70  # measured
config.training.student_baseline = 0.40  # measured

expected = config.training.expected_saturation_accuracy()
print(f"Expected saturation: {expected:.1%}")  # ~65%

# Full analysis
print(config.training.gap_analysis())
\end{lstlisting}

\section{Adaptive Features (v1.2)}

\subsection{Adaptive Complexity Ratio (Figure 7)}

The paper shows that complex data can hurt performance when the teacher model is weak
(LEXam case). The adaptive complexity feature automatically reduces the complexity
ratio when teacher accuracy is below 60\%:

\begin{lstlisting}[language=Python,caption={Adaptive complexity ratio}]
# In simula_config.py GenerationConfig
def get_adaptive_complexity_ratio(self, teacher_accuracy: float | None) -> float:
    """Reduce complexity for weak teachers (Figure 7)."""
    if not self.adaptive_complexity or teacher_accuracy is None:
        return self.complexity_ratio
    
    if teacher_accuracy < 0.6:
        return 0.2  # Significantly reduce
    elif teacher_accuracy < 0.7:
        return min(self.complexity_ratio, 0.35)
    return self.complexity_ratio
\end{lstlisting}

\subsection{Adaptive Critic Thresholding (Figure 3b)}

Critic effectiveness degrades with complexity. The adaptive threshold increases
the acceptance bar for high-complexity examples:

\begin{lstlisting}[language=Python,caption={Adaptive critic threshold}]
def get_adaptive_critic_threshold(self, complexity_score: float | None) -> float:
    """Adjust threshold based on complexity (Figure 3b)."""
    if complexity_score is None:
        return self.critic_threshold
    if complexity_score > 0.7:
        return min(0.7, self.critic_threshold + 0.2)  # More conservative
    elif complexity_score < 0.3:
        return max(0.3, self.critic_threshold - 0.2)  # More permissive
    return self.critic_threshold
\end{lstlisting}

\subsection{Automatic Coverage Gap Filling}

When coverage evaluation identifies under-represented taxonomy nodes, the system
can automatically generate additional examples to fill gaps:

\begin{lstlisting}[language=Python,caption={Coverage gap filling}]
from src.training.pipeline import fill_coverage_gaps, identify_coverage_gaps

# Identify gaps
gaps = identify_coverage_gaps(coverage_reports, threshold=0.5)

# Fill gaps automatically
additional = await fill_coverage_gaps(
    gaps=gaps,
    generator=generator,
    examples_per_node=10,
)
examples.extend(additional)
\end{lstlisting}

\subsection{Diversity Optimizer}

The diversity optimizer adjusts sampling weights to balance global and local diversity:

\begin{lstlisting}[language=Python,caption={Diversity optimizer usage}]
from src.training.pipeline import SimulaDiversityOptimizer, DiversityTargets

optimizer = SimulaDiversityOptimizer(
    targets=DiversityTargets(
        target_global=0.52,  # From Figure 4
        target_local=0.20,
    ),
)

# Analyze current diversity and get recommendations
adjustment = optimizer.analyze(diversity_report)
print(f"Global gap: {adjustment.global_diversity_gap:.3f}")
print(f"Local gap:  {adjustment.local_diversity_gap:.3f}")

# Update weights for next generation batch
new_weights = optimizer.update_weights(diversity_report, coverage_by_taxonomy)
\end{lstlisting}

\subsection{Taxonomy Quality Gate (Section B.1, Table 2)}

Generated taxonomies can be validated against expert references:

\begin{lstlisting}[language=Python,caption={Taxonomy quality gate}]
from src.training.pipeline import SimulaTaxonomyEvaluator, TaxonomyQualityError

evaluator = SimulaTaxonomyEvaluator(llm_client)

# Compare against expert taxonomy
report = await evaluator.compare(generated_taxonomy, "data/expert/taxonomy.json")
print(f"Completeness: {report.completeness:.1%}")  # Should be > 70%
print(f"Soundness: {report.soundness:.1%}")        # Should be > 90%

# Validate thresholds (raises TaxonomyQualityError if fails)
evaluator.validate_taxonomy(generated_taxonomy, report=report)
\end{lstlisting}

\section{Evaluation Environment Variables}

\begin{table}[H]
\centering
\caption{Evaluation Framework Environment Variables}
\label{tab:eval-env}
\begin{tabularx}{\textwidth}{lp{2cm}X}
\toprule
\textbf{Variable} & \textbf{Default} & \textbf{Description} \\
\midrule
\texttt{SIMULA\_COMPLEXITY\_BATCH\_SIZE} & 5 & BS parameter (Algorithm 3) \\
\texttt{SIMULA\_COMPLEXITY\_N\_SAMPLES} & 5 & N parameter (Algorithm 3) \\
\texttt{SIMULA\_COMPLEXITY\_ENABLE\_ELO} & true & Enable Elo calibration \\
\texttt{SIMULA\_COVERAGE\_MAX\_LEVEL} & 4 & Max taxonomy level \\
\texttt{SIMULA\_COVERAGE\_SAMPLE\_SIZE} & 1000 & Coverage sample size \\
\texttt{SIMULA\_DIVERSITY\_EMBEDDING\_MODEL} & MiniLM & Embedding model (or Gecko) \\
\texttt{SIMULA\_DIVERSITY\_KNN} & 10 & k-NN for local diversity \\
\texttt{SIMULA\_CRITIC\_VALIDATION\_SPLIT} & 0.05 & Validation fraction \\
\texttt{SIMULA\_TEACHER\_MODEL} & Qwen3.5-35B & Teacher model name \\
\texttt{SIMULA\_STUDENT\_MODEL} & Gemma-3-4B & Student model name \\
\texttt{SIMULA\_SATURATION\_RATIO} & 0.83 & Gap saturation ratio \\
\midrule
\multicolumn{3}{l}{\textbf{v1.2 Adaptive Features}} \\
\midrule
\texttt{SIMULA\_ADAPTIVE\_COMPLEXITY} & true & Auto-reduce for weak teacher \\
\texttt{SIMULA\_CRITIC\_ADAPTIVE\_THRESHOLD} & true & Complexity-aware rejection \\
\texttt{SIMULA\_CRITIC\_THRESHOLD} & 0.5 & Base critic threshold \\
\texttt{SIMULA\_TAXONOMY\_MIN\_COMPLETENESS} & 0.7 & Minimum completeness \\
\texttt{SIMULA\_TAXONOMY\_MIN\_SOUNDNESS} & 0.9 & Minimum soundness \\
\texttt{SIMULA\_DIVERSITY\_TARGET\_GLOBAL} & 0.52 & Target global diversity \\
\texttt{SIMULA\_DIVERSITY\_TARGET\_LOCAL} & 0.20 & Target local diversity \\
\texttt{SIMULA\_DIVERSITY\_TOLERANCE} & 0.05 & Diversity tolerance \\
\texttt{SIMULA\_DIVERSITY\_LEARNING\_RATE} & 0.1 & Weight adjustment rate \\
\bottomrule
\end{tabularx}
\end{table}

% =============================================================================
% WORKED EXAMPLE: END-TO-END WALKTHROUGH
% =============================================================================
\newpage
\chapter*{Worked Example: Trial Balance Domain}
\addcontentsline{toc}{section}{Worked Example}

This section demonstrates an end-to-end walkthrough of the Simula pipeline for the
Trial Balance domain, showing actual inputs, intermediate outputs, and final results.

\section*{Step 1: Schema Extraction}

\textbf{Input:} HANA schemas \texttt{PAL\_STORE.TB\_VARIANCES}, \texttt{PAL\_STORE.GL\_BALANCES}

\begin{lstlisting}[language=Python,caption={Schema extraction output (schemas.json)}]
{
  "schemas": {
    "PAL_STORE": {
      "tables": {
        "TB_VARIANCES": {
          "columns": [
            {"name": "VARIANCE_ID", "type": "INTEGER", "nullable": false},
            {"name": "ACCOUNT_CODE", "type": "NVARCHAR(20)", "nullable": false},
            {"name": "PERIOD", "type": "NVARCHAR(7)", "nullable": false},
            {"name": "ACTUAL_AMOUNT", "type": "DECIMAL(18,2)", "nullable": true},
            {"name": "BUDGET_AMOUNT", "type": "DECIMAL(18,2)", "nullable": true},
            {"name": "VARIANCE_PCT", "type": "DECIMAL(5,2)", "nullable": true},
            {"name": "DEPARTMENT", "type": "NVARCHAR(50)", "nullable": true}
          ]
        },
        "GL_BALANCES": {
          "columns": [
            {"name": "ACCOUNT_CODE", "type": "NVARCHAR(20)", "nullable": false},
            {"name": "BALANCE", "type": "DECIMAL(18,2)", "nullable": false},
            {"name": "BALANCE_DATE", "type": "DATE", "nullable": false}
          ]
        }
      }
    }
  },
  "extracted_at": "2026-04-18T08:00:00Z"
}
\end{lstlisting}

\section*{Step 2: Factor Identification}

\textbf{LLM Output:} 4 factors identified for Trial Balance domain

\begin{lstlisting}[language=Python,caption={Factors (factors.json)}]
{
  "factors": [
    {"name": "query_type", "description": "Type of SQL operation"},
    {"name": "time_granularity", "description": "Time period level"},
    {"name": "aggregation_scope", "description": "Level of aggregation"},
    {"name": "variance_analysis", "description": "Type of variance check"}
  ]
}
\end{lstlisting}

\section*{Step 3: Taxonomy Generation}

\textbf{Parameters:} depth=3, best\_of\_n=5

\begin{table}[H]
\centering
\caption{Generated Taxonomy for ``query\_type'' Factor}
\begin{tabular}{llll}
\toprule
\textbf{Level 0} & \textbf{Level 1} & \textbf{Level 2} & \textbf{Nodes} \\
\midrule
query\_type & aggregation & sum, count, avg, max, min & 5 \\
            & filtering & by\_threshold, by\_period, by\_department & 3 \\
            & comparison & MoM, YoY, budget\_vs\_actual & 3 \\
            & ranking & top\_N, bottom\_N, percentile & 3 \\
\midrule
\multicolumn{3}{l}{\textbf{Total nodes (query\_type)}} & 15 \\
\bottomrule
\end{tabular}
\end{table}

\textbf{Total across all 4 factors:} 52 leaf nodes

\section*{Step 4: Mix Sampling and Meta-Prompt Generation}

\textbf{Sample mix:} \texttt{[aggregation/sum, MoM, by\_department]}

\begin{lstlisting}[language=Python,caption={Meta-prompt generated}]
"Generate a natural language question asking for the sum of 
variances by department, comparing month-over-month changes, 
for the Trial Balance review process."
\end{lstlisting}

\section*{Step 5: Example Generation}

\begin{lstlisting}[language=Python,caption={Generated training example}]
{
  "id": "tb-sim-001",
  "question": "What is the total variance amount by department 
               for November vs October 2025?",
  "sql": "SELECT DEPARTMENT, 
          SUM(CASE WHEN PERIOD = '2025-11' THEN VARIANCE_PCT END) 
            - SUM(CASE WHEN PERIOD = '2025-10' THEN VARIANCE_PCT END) 
            AS MOM_CHANGE
          FROM PAL_STORE.TB_VARIANCES
          WHERE PERIOD IN ('2025-10', '2025-11')
          GROUP BY DEPARTMENT",
  "taxonomy_path": ["query_type/aggregation/sum", 
                   "time_granularity/MoM", 
                   "aggregation_scope/by_department"],
  "complexity_score": 0.65,
  "is_complexified": false,
  "critic_passed": true
}
\end{lstlisting}

\section*{Step 6: Evaluation Metrics}

After generating 10,247 examples for Trial Balance in the reference
run of 2025-11-15 (commit \texttt{b8d40a1},
tag \texttt{simula-v1.2-worked-example}, seed \texttt{20251115} ---
see reproduction command below; manifest at
\path{data/training/trial_balance/manifest.json}):

\begin{table}[H]
\centering
\caption{Evaluation results for Trial Balance dataset
(reference run \texttt{simula-v1.2-worked-example}, 2025-11-15).}
\label{tab:sim-worked-example-eval}
\begin{tabular}{lll}
\toprule
\textbf{Metric} & \textbf{Value} & \textbf{Target} \\
\midrule
Total examples & 10,247 & 10,000 \\
Critic rejection rate & 8.2\% & $<$15\% \\
Global diversity & 0.51 & 0.52 \\
Local diversity & 0.22 & 0.20 \\
Coverage (Level 1) & 100\% & 100\% \\
Coverage (Level 2) & 92\% & $>$90\% \\
Coverage (Level 3) & 78\% & $>$70\% \\
Mean Elo score & 1,512 & -- \\
High complexity split & 38\% & 40\% \\
\bottomrule
\end{tabular}
\end{table}

\subsubsection{Reproducibility.}
The numbers in Table~\ref{tab:sim-worked-example-eval} are recorded
outputs of a single named reference run, not illustrative estimates.
Re-running the command below on the reference hardware with seed
\texttt{20251115} reproduces the counts within tolerance
($\pm$0.5\,\% on the critic rejection rate and $\pm$0.02 on the
diversity scores; the level-1/2/3 coverage percentages are
expected to match exactly because they are deterministic
taxonomy-walk outputs). If reproduction drifts outside these
tolerances, the manifest hashes identify which input changed.

\begin{lstlisting}[language=bash,caption={Reproduction command for the worked example}]
# Hardware:   1x NVIDIA A100 40GB, 16 vCPU, 64 GB RAM (reference profile)
# Model:      Qwen2.5-32B-Instruct served via vLLM 0.6.3
# Commit:     b8d40a1 (tag: simula-v1.2-worked-example)
# Wall time:  ~5h 40m end-to-end

python -m src.training.pipeline.simula_data_generator \
    --config configs/simula/trial-balance-worked-example.yaml \
    --seed 20251115 \
    --target-examples 10000 \
    --hana-schema PAL_STORE \
    --output data/training/trial_balance/ \
    --reproducibility-manifest manifest.json
\end{lstlisting}

\noindent
The reproducibility manifest (\texttt{manifest.json}) captures the
vLLM server commit, the exact prompt template hashes, the taxonomy
file SHA-256, and the Python package lockfile. Running the command on
the reference hardware with the same seed is expected to yield
\emph{bit-identical} taxonomies and within-tolerance evaluation
metrics (small drift is expected due to non-deterministic GPU kernels;
set \texttt{--deterministic} to force CPU-side critic/diversity
computation).

\section*{Step 7: Output Files}

\begin{lstlisting}[language=bash,caption={Final output artifacts}]
data/training/trial_balance/
  training_data.jsonl        # 10,247 examples (2.4 MB)
  schemas.json               # Extracted HANA schemas
  taxonomies.json            # 4 taxonomies, 52 leaf nodes
  complexity_elo_scores.json # Elo scores for all examples
  coverage_report.json       # Level ratio coverage
  diversity_metrics.json     # Global=0.51, Local=0.22
  critic_calibration.json    # p(y)=0.87, p(corrupt)=0.12
  generation_stats.json      # Timing, counts, rates
  rejected_examples.jsonl    # 923 rejected by critic
\end{lstlisting}

% =============================================================================
% CHAPTER 8: REFERENCES
% =============================================================================
\newpage
\chapter{References}
\label{chap:references}

\section{Source Documents}

\begin{table}[H]
\centering
\caption{Source Documents}
\begin{tabularx}{\textwidth}{lX}
\toprule
\textbf{Document} & \textbf{Path} \\
\midrule
Simula Research Paper & \texttt{docs/6171\_Reasoning\_Driven\_Syntheti.pdf} \\
Trial Balance Spec & \texttt{docs/technical-reports/04-trial-balance-ai-specification.tex} \\
Arabic AP Spec & \texttt{docs/arabic/structured/} (external) \\
\bottomrule
\end{tabularx}
\end{table}

\section{Implementation Files}

\begin{table}[H]
\centering
\caption{Implementation Files}
\begin{tabularx}{\textwidth}{lX}
\toprule
\textbf{File} & \textbf{Lines} \\
\midrule
\texttt{src/training/pipeline/simula\_config.py} & Configuration dataclasses \\
\texttt{src/training/pipeline/simula\_llm\_client.py} & LLM client with Best-of-N, double-critic \\
\texttt{src/training/pipeline/hana\_schema\_extractor.py} & HANA MCP extraction \\
\texttt{src/training/pipeline/simula\_taxonomy\_builder.py} & Algorithm 1 implementation \\
\texttt{src/training/pipeline/simula\_data\_generator.py} & Algorithm 2 implementation \\
\texttt{src/training/pipeline/simula\_complexity\_calibrator.py} & Algorithm 3 (Elo scoring) \\
\texttt{src/training/pipeline/simula\_coverage\_evaluator.py} & Level Ratio Coverage \\
\texttt{src/training/pipeline/simula\_diversity\_analyzer.py} & Global + local diversity \\
\texttt{src/training/pipeline/simula\_critic\_validator.py} & Critic calibration \\
\texttt{src/training/scripts/run\_hana\_pipeline.py} & CLI entry point \\
\texttt{src/training/scripts/split\_by\_complexity.py} & Complexity split (Figure 7) \\
\bottomrule
\end{tabularx}
\end{table}

\section{Tools}

\begin{table}[H]
\centering
\caption{Tools and Dependencies}
\begin{tabular}{lll}
\toprule
\textbf{Tool} & \textbf{Version} & \textbf{Purpose} \\
\midrule
vLLM TurboQuant & 0.8.x & LLM inference backend \\
AI Core PAL & 1.x & HANA MCP server \\
Python & 3.11+ & Runtime \\
asyncio & stdlib & Async orchestration \\
httpx & 0.27+ & HTTP client \\
\bottomrule
\end{tabular}
\end{table}
